% Generated by tools/edn_latex.py from reports/edn.md.
% Do not edit: change reports/edn.md and run `make edn`.
\ednentry{
  id = {EDN-69},
  title = {CodeCarbon measures the fit alone, offline and per process, and its figure is reported as the estimate it is},
  date = {2026-10-05},
  milestone = {M3: Quality Assurance},
  topic = {Energy Efficiency, Experiment Tracking, Evaluation Protocol},
  participants = {@lukas2510, decided by the repository owner},
  decision = {\texttt{train} measures exactly the \texttt{fit\_\allowbreak{}variant} call with CodeCarbon 3.3.1's \texttt{OfflineEmissionsTracker}, in \texttt{process} mode, against the Spanish grid (\texttt{ESP}, 174.05 g CO2eq/kWh) and with \texttt{measure\_\allowbreak{}power\_\allowbreak{}secs: 15}, all three from \texttt{params.\allowbreak{}yaml}'s \texttt{train.\allowbreak{}energy}, which the \texttt{train} stages declare. CodeCarbon's row goes to \texttt{reports/\allowbreak{}emissions/\allowbreak{}<variant>.\allowbreak{}csv} and into the MLflow run that also receives the variant's test metrics, as metrics, with the CPU power method and the machine behind them as parameters. A \texttt{compare-\allowbreak{}energy} stage joins each model, its row and its metrics into \texttt{reports/\allowbreak{}figures/\allowbreak{}energy-\allowbreak{}vs-\allowbreak{}error.\allowbreak{}\{csv,png\}}. On every machine we train on the figure is CodeCarbon's estimate, not a reading: 10 W for the RAM plus the CPU's TDP times the share of its logical CPUs the fit kept busy. We keep that estimate in the pipeline, word every comparison as one of fit times expressed in energy, within one run on one machine, and checked the estimate once against the RAPL energy counter, with read access granted for that check only.},
  rationale = {\emph{Alternatives considered:}
\begin{itemize}[leftmargin=*, topsep=2pt]
\item \textbf{How the figure is obtained, given that it is an estimate.}
\par
\begin{itemize}[leftmargin=*, topsep=2pt]
\item \textbf{Option A: keep CodeCarbon's estimate and word it as what it is.}
\newline Pros: nothing changes on anyone's machine; the same mechanism and settings everywhere, hermetic, with no network call; it is the prescribed tool used as prescribed.
\newline Cons: the energy axis adds no information beyond the fit time it is computed from, and the absolute watt-hours are model output that depends on the machine and on its load.
\item \textbf{Option B: grant read access to RAPL on Lukas's laptop permanently,} with a udev or tmpfiles rule, so that CodeCarbon reads the package counter there.
\newline Pros: a hardware counter rather than a model.
\newline Cons: RAPL is root-only on purpose since kernel 5.10, because unprivileged power readings are a side channel that has been used to extract keys (Platypus, CVE-2020-8694), and this laptop has SGX. RAPL counts the whole package and CodeCarbon does not apportion it to the process, so on a laptop several agents share -- the 1-minute load ranged from 10 to 31 on the day of the measurement -- it would charge each fit for the others. Only Lukas's runs would read it, so a teammate's \texttt{dvc repro} would write a figure of another kind into the same files. And it is a standing change to a personal machine that somebody has to remember.
\item \textbf{Option C (chosen): A, plus a one-off check of the estimate against the counter.} Read access is granted with a \texttt{chmod} that does not survive a reboot, the four fits and an idle baseline are measured with both, and access is revoked again.
\newline Pros: the report can say what the estimate is worth instead of only what it is, at no standing cost to the laptop's security, and the pipeline stays as in A.
\newline Cons: it needs Lukas's \texttt{sudo} and an otherwise idle laptop, and it validates one machine on one day, not every run.
\item \textbf{Option D: compare the ladder on fit seconds and CPU-seconds, and keep CodeCarbon's figure only as NFR-10's log.}
\newline Pros: both are direct measurements.
\newline Cons: it drops the energy comparison issue \#38 and the rubric ask for, and the comparison would show the same ratios under another name.
\end{itemize}
\item \textbf{\texttt{process} against \texttt{machine} mode.} \texttt{machine} charges the fit for the load of the whole machine: on the shared laptop, a process that only slept was charged 4.5 W of CPU in \texttt{machine} mode and 0.3 W in \texttt{process} mode. \texttt{process} reads the fit's own CPU time.
\item \textbf{The online tracker against the offline one with a pinned country.} Per CodeCarbon's source, the online tracker geolocates the machine through a call to geojs, writes its latitude and longitude into the CSV, and queries the AWS, Azure and GCP metadata endpoint, switching to the provider's regional factor when one answers, which on a GitHub runner, an Azure VM, one would. The emissions of an identical fit would then depend on where it ran, a git-tracked file would carry a contributor's location, and the test suite would make network calls. With every socket connection made to fail, the offline tracker made none.
\item \textbf{The demo's \texttt{measure\_\allowbreak{}power\_\allowbreak{}secs=\allowbreak{}1} against CodeCarbon's default of 15.} The power is sampled every second by a separate thread either way; this setting is how often the samples are integrated, and at 1 s that integration races the tracker's own shutdown. Against CodeCarbon's own model, 1 s gave 0.86x to 1.07x on a single-threaded LightGBM fit and up to 4.31x on a busy loop's RAM half, an interval counted several times over; 15 s gave 0.987x to 0.993x on the same fit.
\end{itemize}
\par
\emph{Why:} Issue \#38 was written on the assumption that CodeCarbon measures the energy of a fit, and the first thing done was to test that on the laptop the ladder is trained on, before any number was quoted. It does not. Since kernel 5.10 RAPL is root-only, so CodeCarbon falls back to its \texttt{cpu\_\allowbreak{}load} model, with the i5-10210U's TDP of 15 W from its own table and a constant 10 W for the RAM, its floor of two DIMMs for any x86 machine. The energy of a fit is then 10 W times its wall time plus 15 W / 8 logical CPUs times its CPU time, and at \texttt{train.\allowbreak{}num\_\allowbreak{}threads: 1} that is its duration times 11.75 to 11.80 W, to within 1\,\% over repeated fits. The RAM constant is 83\,\% of a LightGBM fit's figure. On a GitHub runner, an AMD EPYC 7763 with four vCPUs and no powercap interface at all, the same model charges 80 W per second, because it divides the chip's 280 W TDP by the four vCPUs the VM sees. So the figure is a fit time multiplied by a machine-specific constant, and it moves with the machine's load as well as with the work: the same \texttt{lgbm-\allowbreak{}basic} fit took 33.0 s on a laptop shared with other jobs and 8.6 s on a quiet one, and was charged accordingly.
\par
That makes the comparison the issue asked for weaker than it reads, not wrong. Between two variants fitted in one run on one machine the ratio of their energies is the ratio of their fit times, which is a fair measure of their relative cost. What it cannot be is a statement in watt-hours that holds across machines, or evidence that a variant uses energy differently from how it uses time. So the pipeline keeps the estimate (option A), every place that shows the figure says what it is -- the CSV, the \texttt{codecarbon.\allowbreak{}cpu\_\allowbreak{}power\_\allowbreak{}method} parameter of every run, the subtitle of the figure, the model card and the report -- and the one check that can say how far the estimate is from the hardware was made once (option C) rather than built into every run (option B). NFR-10's specification says the same since this decision, because ``CodeCarbon measures'' was not true of any machine we have.
\par
The RAPL check, \texttt{reports/\allowbreak{}analysis/\allowbreak{}rapl\_\allowbreak{}validation.\allowbreak{}py}, measured an idle baseline and three fits of each variant on an idle laptop with the package, core, uncore, DRAM and platform counters, CodeCarbon's forced \texttt{cpu\_\allowbreak{}load} estimate and CodeCarbon's own RAPL figure, all around the same fit. \textbf{[pending RAPL check: idle draw, estimate against the counter per variant, whether the ranking and the ratios hold.]}
\par
The rest of the design follows from treating the figure as evidence that has to be checkable. Only the fit is inside the tracker, as the demo does, which \texttt{tests/\allowbreak{}test\_\allowbreak{}model.\allowbreak{}py::test\_\allowbreak{}only\_\allowbreak{}the\_\allowbreak{}fit\_\allowbreak{}is\_\allowbreak{}measured} asserts by making the read and the write slow. Every CodeCarbon argument is passed explicitly, because it reads \texttt{\textasciitilde{}/\allowbreak{}.\allowbreak{}codecarbon.\allowbreak{}config}, \texttt{.\allowbreak{}/\allowbreak{}.\allowbreak{}codecarbon.\allowbreak{}config} and any \texttt{CODECARBON\_\allowbreak{}*} variable for every argument it is not given, so a contributor's configuration could switch on its API output or force a CPU power. CodeCarbon swallows every exception in its constructor, \texttt{start} and \texttt{stop}, so the module reads its row back by CodeCarbon's run id and by column name and fails the stage when it is missing, and refuses a country CodeCarbon would silently measure against the world average. The row is joined to its model by that run id, which \texttt{model.\allowbreak{}json} records, and \texttt{fit\_\allowbreak{}cpu\_\allowbreak{}seconds} is logged beside the energy because it is a measurement and the CPU half is proportional to it. DVC deletes a stage's outputs before running it, so the CSV CodeCarbon appends to holds exactly one row after \texttt{dvc repro} (measured: one row after three \texttt{dvc repro -\allowbreak{}-\allowbreak{}force}, three with \texttt{persist: true}).
\par
What the comparison concluded, on the ladder of \textbf{[pending final run: commit]}: \textbf{[pending final run: per-variant fit time, energy, emissions and MdAPE, and whether \texttt{lgbm-\allowbreak{}extended} earns its extra training energy]}. The whole ladder is \textbf{[pending final run]} of fitting, against NFR-10's 15-minute budget.},
  ai = {Information seeking, Alternative generation, Alternative assessment, Recommendation, Solution generation},
  response = {Accepted},
  assessment = {AI found the limitation the decision is about, by reading CodeCarbon's source and then measuring it on the laptop and on a GitHub runner, before writing the module; the issue, the course demo and the report draft all assumed a measurement. It laid out the four options with the security and the shared-machine arguments against B, recommended A with C, and Lukas chose that. It also chose the tracker's settings on measurements rather than on the demo: \texttt{process} mode, the offline tracker, and 15 s rather than the demo's 1 s once it was shown to race the tracker's shutdown. It did not grant itself RAPL access by other means, although Lukas's \texttt{docker} group would have allowed it, because that would have bypassed the permission boundary this decision is partly about.},
  aievidence = {Claude Code session on 2026-10-05 implementing issue \#38: probes of CodeCarbon 3.3.1 with every socket blocked, on four workloads in both modes and at both integration intervals, a DVC repository testing the append, and a temporary CI step on pull request \#68 that printed what CodeCarbon does on a GitHub runner; then the four options above with a recommendation, which Lukas decided as option C.},
  evidence = {\href{https://github.com/mlops-2627q1-mds-upc/MLOPS_Recommenditos/blob/main/recommenditos/modeling/energy.py}{\texttt{recommenditos/\allowbreak{}modeling/\allowbreak{}energy.\allowbreak{}py}}, whose docstring states what the figure is, and \href{https://github.com/mlops-2627q1-mds-upc/MLOPS_Recommenditos/blob/main/recommenditos/modeling/compare_energy.py}{\texttt{recommenditos/\allowbreak{}modeling/\allowbreak{}compare\_\allowbreak{}energy.\allowbreak{}py}}; \texttt{train.\allowbreak{}energy} in \href{https://github.com/mlops-2627q1-mds-upc/MLOPS_Recommenditos/blob/main/params.yaml}{\texttt{params.\allowbreak{}yaml}}; \href{https://github.com/mlops-2627q1-mds-upc/MLOPS_Recommenditos/blob/main/tests/test_energy.py}{\texttt{tests/\allowbreak{}test\_\allowbreak{}energy.\allowbreak{}py}} and \texttt{tests/\allowbreak{}test\_\allowbreak{}model.\allowbreak{}py::test\_\allowbreak{}every\_\allowbreak{}training\_\allowbreak{}run\_\allowbreak{}records\_\allowbreak{}its\_\allowbreak{}emissions\_\allowbreak{}beside\_\allowbreak{}its\_\allowbreak{}accuracy} and \texttt{::test\_\allowbreak{}only\_\allowbreak{}the\_\allowbreak{}fit\_\allowbreak{}is\_\allowbreak{}measured}; \href{https://github.com/mlops-2627q1-mds-upc/MLOPS_Recommenditos/blob/main/reports/analysis/codecarbon_validity.py}{\texttt{reports/\allowbreak{}analysis/\allowbreak{}codecarbon\_\allowbreak{}validity.\allowbreak{}py}} with its results, \href{https://github.com/mlops-2627q1-mds-upc/MLOPS_Recommenditos/blob/main/reports/analysis/codecarbon_github_runner.txt}{\texttt{codecarbon\_\allowbreak{}github\_\allowbreak{}runner.\allowbreak{}txt}} and \href{https://github.com/mlops-2627q1-mds-upc/MLOPS_Recommenditos/blob/main/reports/analysis/rapl_validation.py}{\texttt{rapl\_\allowbreak{}validation.\allowbreak{}py}} with its results; \texttt{reports/\allowbreak{}figures/\allowbreak{}energy-\allowbreak{}vs-\allowbreak{}error.\allowbreak{}png}; \href{https://github.com/mlops-2627q1-mds-upc/MLOPS_Recommenditos/blob/main/docs/docs/specification.md}{specification} NFR-10; \href{https://github.com/mlops-2627q1-mds-upc/MLOPS_Recommenditos/blob/main/docs/docs/model-card.md}{model card}, Environmental Impact; \href{https://github.com/mlops-2627q1-mds-upc/MLOPS_Recommenditos/issues/38}{issue \#38}, \href{https://github.com/mlops-2627q1-mds-upc/MLOPS_Recommenditos/pull/68}{pull request \#68}.},
}
